%-------------------------------------------------------------------------------
%  Shared preamble for cv-en.tex and cv-vi.tex
%  Edit the layout HERE, never in the content files — this is what keeps the
%  English and Vietnamese versions looking identical.
%
%  ENGINE: XeLaTeX (Overleaf: Menu → Compiler → XeLaTeX) or Tectonic.
%  NOT pdfLaTeX — fontspec needs a Unicode engine, and cv-vi.tex needs one for
%  the Vietnamese diacritics anyway.
%
%  LOOK: modelled on the Jake Gutierrez resume template Nhi supplied —
%  black and white, small-caps section headings over a full-width rule,
%  skills as "Label: values" lines, bold project title left / date right.
%
%  TYPOGRAPHY
%  * Body is TeX Gyre Termes (a Times cut): the classic CV serif, familiar to
%    every recruiter, and far sturdier than Computer Modern (the LaTeX default),
%    whose thin strokes and small x-height are punishing at small sizes.
%  * DO NOT swap in another font without checking Vietnamese first. XCharter was
%    tried here and silently dropped 48 Vietnamese glyphs — ả, ẩ, ủ, ộ, ể and
%    friends just vanished from cv-vi.pdf while the build still reported success.
%    To verify a candidate font, the build log must show zero hits for:
%      grep -c "Missing character" <build log>      # must be 0
%  * Fonts are loaded BY FILENAME, not family name. Loading "TeX Gyre Termes" by
%    name fails on Tectonic (no system font database); the .otf filename resolves
%    on Overleaf, TeX Live and Tectonic alike.
%-------------------------------------------------------------------------------

\usepackage[
  a4paper,
  top=1.1cm, bottom=1.1cm, left=1.4cm, right=1.4cm,
  footskip=0.4cm
]{geometry}

\usepackage{fontspec}
\usepackage{fontawesome}
\usepackage{titlesec}
\usepackage{enumitem}
\usepackage{array}
\usepackage{tabularx}
\usepackage{ragged2e}
\usepackage{xcolor}
\usepackage[hidelinks]{hyperref}
\usepackage{fancyhdr}

% ─── Fonts ──────────────────────────────────────────────────────────────────
\setmainfont{texgyretermes-regular.otf}[
  BoldFont       = texgyretermes-bold.otf,
  ItalicFont     = texgyretermes-italic.otf,
  BoldItalicFont = texgyretermes-bolditalic.otf,
  Ligatures      = TeX,
]

% ─── Colour: black on white, like the original. Links only are blue. ────────
\definecolor{link}{HTML}{0B3D91}

\hypersetup{colorlinks=true, urlcolor=link}
\urlstyle{same}

\pagestyle{fancy}
\fancyhf{}
\renewcommand{\headrulewidth}{0pt}
\renewcommand{\footrulewidth}{0pt}

\setlength{\parindent}{0pt}
\setlength{\tabcolsep}{0in}
\linespread{1.02}
\raggedright
\raggedbottom

% ─── Sections: small caps over a full-width rule (the Jake look) ────────────
\titleformat{\section}
  {\scshape\large}
  {}{0em}{}
  [\vspace{-6pt}\rule{\linewidth}{0.8pt}\vspace{-3pt}]
\titlespacing*{\section}{0pt}{8pt}{4pt}

% ─── Entry macros ───────────────────────────────────────────────────────────
% \cvHeading{Title}{Right-aligned date}{Role / subtitle}
\newcommand{\cvHeading}[3]{%
  \vspace{3pt}
  \begin{tabular*}{\linewidth}{@{}l@{\extracolsep{\fill}}r@{}}
    \textbf{#1} & \textit{#2} \\
  \end{tabular*}%
  \vspace{-2pt}\\
  {\small\itshape #3}%
  \vspace{1pt}
}

% \cvSchool{School}{Location}{Degree line}{Dates}
\newcommand{\cvSchool}[4]{%
  \vspace{3pt}
  \begin{tabular*}{\linewidth}{@{}l@{\extracolsep{\fill}}r@{}}
    \textbf{#1} & \textit{#2} \\
  \end{tabular*}%
  \vspace{-2pt}\\
  \begin{tabular*}{\linewidth}{@{}l@{\extracolsep{\fill}}r@{}}
    {\small\itshape #3} & {\small\itshape #4} \\
  \end{tabular*}%
  \vspace{-6pt}
}

% The italic "Problem:" line that opens each project.
% \noindent is required — \justifying restores \parindent, which would otherwise
% indent the first line and break the left edge of the column.
\newcommand{\cvProblem}[1]{%
  \vspace{2pt}
  {\small\itshape\justifying\noindent #1\par}%
  \vspace{1pt}
}

% Trailing "Key Tech / Công nghệ" line.
\newcommand{\cvMeta}[2]{%
  \vspace{2pt}
  {\small\textit{#1:} #2\par}%
}

% Project bullets.
\newenvironment{cvBullets}
  {\begin{itemize}[
     leftmargin = 1.3em,
     topsep     = 2pt,
     itemsep    = 1.5pt,
     parsep     = 0pt,
     label      = $\bullet$,
   ]\small}
  {\end{itemize}\vspace{1pt}}

% Flat, label-less list — used for Summary, Skills and Certifications, matching
% the original template's "\textbf{Label}{: values}" style.
\newenvironment{cvPlain}
  {\begin{itemize}[leftmargin=0pt, topsep=2pt, itemsep=3pt, parsep=0pt, label={}]\small}
  {\end{itemize}}
